\chapter{投资银行：资本筹集与资本市场中介}
\label{chap:Investment_Banks_Capital_Raising_and_Capital_Market_Intermediation}

\begin{openingcase}{直接融资为什么需要投资银行？}

这家发展迅速的公司计划通过向公众投资者出售新股来筹集约60亿元人民币的资金。该公司的财务报表已经过审计，相关信息可以在招股书中找到。投资者最终将持有该公司发行的股票。

然而，该公司并不会独自与投资者进行接触。它会聘请一家投资银行，并支付相应的费用以获得其服务。这家投资银行会深度参与项目的筹备、定价以及发行过程，并且可能会承担与本次发行相关的一部分风险。

\end{openingcase}

\begin{tcolorbox}[
  enhanced,
  colback=FinanceBlue!6,
  colframe=FinanceBlue,
  boxrule=0.6pt,
  sharp corners,
  left=8pt,
  right=8pt,
  top=7pt,
  bottom=7pt,
  before skip=1em,
  after skip=1.6em,
  title=\textbf{核心问题},
  coltitle=white,
  colbacktitle=FinanceBlue,
  fonttitle=\bfseries,
  attach boxed title to top left={xshift=6pt,yshift=-2mm},
  boxed title style={
    boxrule=0pt,
    sharp corners,
    left=5pt,
    right=5pt,
    top=2pt,
    bottom=2pt
  }
]
\begin{enumerate}[
  leftmargin=*,
  label=\textcolor{FinanceBlue}{\textbf{\arabic*.}},
  itemsep=0.55em,
  topsep=0.2em
]

\item 为什么在证券已经建立了发行人与投资者之间的直接请求权的情况下，发行人仍然使用投资银行呢？

\item 投资银行如何利用信息、声誉、估值、分销网络以及资产负债表能力来支持证券的发行、交易和战略交易呢？

\end{enumerate}
\end{tcolorbox}

\section{资本市场中介与认证}

投资银行帮助发行人获取资本市场资金，而不会取代投资者持有的证券。他们的角色与商业银行的资产负债表中介业务不同，后者通常通过存款和其他银行负债来为借款人的请求权提供资金支持。

\subsection{发行人、投资者与资本市场准入}

发行方进入资本市场时有着不同的融资需求。成长中的公司可以通过发行权益来为投资提供资金，而无需承担固定的偿还义务；而公司则可以通过发行债券来获得短期融资、重新融资现有债务或分散其资金来源。政府和金融机构也根据不同的融资目标和合同结构来发行证券。

进入资本市场需要将融资需求转化为投资者能够评估和购买的证券。发行人必须提供相关信息，明确请求权的条款，吸引潜在投资者，并确定能够吸引足够投资者需求的条款。法律、信息披露和交易成本也决定了发行的规模和设计在经济上的重要性。

投资银行帮助连接市场的这两端。他们就发行的结构和时机提供建议，组织筹备和分配过程，并收集有关投资者需求的信息。这种协调可以减少搜索和交易成本，同时提高发行人与投资者之间的信息流通。图~\ref{fig:ch13-direct-finance-intermediation}总结了这些关系。

\begin{figure}[H]

\centering

\begin{tikzpicture}[
  >=Latex,
  box/.style={
    draw,
    rounded corners,
    align=center,
    minimum width=1.8cm,
    minimum height=0.9cm,
    inner sep=5pt
  },
  service/.style={Latex-Latex, line width=0.85pt, FinanceBlue},
  claim/.style={Latex-Latex, dashed, line width=0.75pt, draw=gray!75}
]

\node[box] (issuer) at (-4.6,0) {发行方};
\node[box] (ib) at (0.5,0) {投资银行};
\node[box] (investors) at (5.2,0) {投资者};

\draw[service] (issuer) --
  node[midway, above=2pt, font=\small,
  text=FinanceBlue, align=center, text width=2.8cm]
  {咨询与交易事项\\协调处理。}
  (ib);

% Investment bank to investors: distribution
\draw[->, line width=0.85pt, FinanceBlue]
  ([yshift=0.14cm]ib.east) --
  node[midway, above=3pt, font=\small,
       text=FinanceBlue, align=center]
  {证券\\分配}
  ([yshift=0.14cm]investors.west);

% Investors to investment bank: demand information
\draw[->, dashed, line width=0.75pt, FinanceBlue]
  ([yshift=-0.14cm]investors.west) --
  node[midway, below=1pt, font=\small,
       align=center]
  { 保留LaTeX命令、公式、交叉引用和数字。\\信息}
  ([yshift=-0.14cm]ib.east);

% Funds: Issuer -> Investors
\coordinate (funds-start) at
  ($(investors.south west)!0.25!(investors.south east)$);

\coordinate (funds-end) at
  ($(issuer.south west)!0.75!(issuer.south east)$);

\coordinate (funds-level) at
  ($(funds-start)!0.5!(funds-end)+(0,-0.75cm)$);

\draw[->, dashed, line width=0.75pt, FinanceBlue]
  (funds-start)
  -- ++(0,-0.75cm)
  -| (funds-end);

\node[below=2pt, font=\small] at (funds-level) {资金};

% Funds: Investors -> Issuer
\coordinate (funds-start) at
  ($(investors.south west)!0.75!(investors.south east)$);

\coordinate (funds-end) at
  ($(issuer.south west)!0.25!(issuer.south east)$);

\coordinate (funds-level) at
  ($(funds-start)!0.5!(funds-end)+(0,-1.5cm)$);

\draw[->, dashed, line width=0.75pt, FinanceBlue]
  (funds-end)
  -- ++(0,-1.5cm)
  -| (funds-start);

\node[below=2pt, font=\small] at (funds-level) {证券};
  
\end{tikzpicture}

\caption{投资银行对资本市场准入的支持}

\label{fig:ch13-direct-finance-intermediation}

\end{figure}


\subsection[尽职调查、信息认证与声誉]{尽职调查、信息认证与声誉}

证券的发行依赖于对发行人、证券以及交易风险的可靠披露。尽职调查通过审查可能影响披露、估值或执行的信息来支持这一过程。尽职调查可能包括对财务报表、商业计划、重要合同、资本结构、所有权、诉讼以及特定交易风险等方面的审查。法律顾问、审计师以及其他专家通常会与投资银行一起参与这一过程。

声誉会影响这项调查背后的激励因素。发行人可能只是偶尔进入资本市场，而投资银行则通常会参与多次交易。因此，如果筛选机制薄弱或信息披露不可靠，就会降低投资者的信心，从而使投资银行在未来的交易中面临更大的困难或更高的成本。

这种声誉风险使投资银行的参与可能发挥认证作用。投资者若无法以合理成本重复同样的调查，就可能认为其参与具有信息价值，因为中介未来业务部分取决于维持市场信誉。


\begin{tcolorbox}[
  enhanced,
  breakable,
  colback=white,
  colframe=FinanceBlue,
  boxrule=0.7pt,
  sharp corners,
  borderline west={2.2pt}{0pt}{FinanceBlue},
  title=\textbf{财务展望：认证并不能保证成功},
  coltitle=white,
  colbacktitle=FinanceBlue,
  fonttitle=\bfseries,
  before skip=1.2em,
  after skip=1.2em
]

认证能够缓解信息摩擦，却不能消除投资风险。投资银行可以调查发行人、审查披露，并通过参与发行使自身声誉承担风险。这些行为可能使其参与本身向投资者传递信息。

但投资银行并不保证发行人未来现金流或证券市场价值。认证影响信息的可信度，无法消除基础请求权的不确定性。

\end{tcolorbox}

\subsection{证券设计与发行人—投资者匹配}

证券设计将发行人的融资需求转化为投资者可以评估和持有的请求权。相关的选择可能包括债务与权益、期限、清偿顺位、支付结构、抵押品、契约限制条款、货币种类以及发行规模等。

这些条款在发行人和投资者之间分配风险和控制权。较长的到期时间可以减少近期再融资的压力，同时增加投资者对利率或信用变化的承受力。抵押品可以加强债权人的保护，同时限制发行人对质押资产的使用。权益则避免了固定的债务偿还，但会改变所有权和请求权。因此，证券的设计涉及权衡取舍，而不是单一的首选结构。

合适的设计还取决于潜在的投资者群体。机构投资者的投资期限、流动性需求、货币风险敞口、监管限制以及允许的资产类别各不相同。即使某项证券在财务上表现良好，但如果其条款不符合潜在买家的要求，那么这项证券可能很难被出售。

投资银行在促成这笔交易的过程中发挥着重要作用。他们利用与投资者的关系、可比的证券、当前市场状况以及预售反馈来评估需求。随着这些信息的不断积累，发行方可以调整债券的规模、期限、定价条款或其他相关因素。因此，证券设计和投资者匹配的过程是同步进行的，直到交易达成能够满足发行方融资需求的同时又能吸引足够多投资者的需求为止。

\section{承销与证券分销}

承销将计划中的证券发行转化为实际的市场交易行为。投资银行负责组织发行过程，与投资者协商证券的定价，并可能承担部分分配风险。

\subsection{承销承诺与承销团安排}

承销协议决定发行人与投资银行如何分担分销风险。在包销安排下，承销商从发行人或出售股份的股东处购买证券，再向投资者出售，因此承担无法按拟定价格完成配售的风险。在尽力销售的代销安排下，投资银行承诺推销证券，但不作相同的购买承诺。

对于具有购买承诺的承销发行，设~$P_o$~为公开发行价格，$P_u$~为承销商向发行人或出售股份的股东支付的价格。每单位证券的承销总价差为

\begin{equation}
g=P_o-P_u,
\label{eq:ch13-underwriting-spread}
\end{equation}
对于$N$种证券而言，按照规定的发行条款，总的毛差价为

\begin{equation}
G=N(P_o-P_u).
\label{eq:ch13-total-underwriting-spread}
\end{equation}

价差补偿承销、分销及相关交易服务，并不等于利润。承销商承担经营和融资成本，投资者需求减弱还可能降低尚未配售证券的实现价值。

大规模发行常组建承销团，结合不同机构的分销能力，并分担承销风险暴露。牵头行协调交易，其他成员承担约定的承销与分销职责。承销团可扩大投资者覆盖，但也需要协调职责、收益分配、客户覆盖与配售。

\begin{example}{一笔包销发行}{firm-commitment-underwriting}

一家公司通过承销方式发行了2000万股新股。公开发行的价格为每股约25元人民币，而承销商则从发行人处以每股约24.25元人民币的价格购得了这些股票。

总差价为

\[
25-24.25=\text{RMB }0.75
\]

每股的价格，而总毛利润差距为

\[
20{,}000{,}000\times0.75=\text{RMB }15\text{ 百万}.
\]

在支付其他费用之前，发行人会先获得收益。

\[
20{,}000{,}000\times24.25=\text{RMB }485\text{ 百万}.
\]

如果在做出承销承诺后，投资者需求减弱，那么部分股票可能需要以低于发行价格的价格出售，或者暂时持有。因此，合同中的价差并不能消除承销商在分配条件方面的风险。

\end{example}


\subsection{首次公开发行与发行定价}

首次公开发行（IPO）是指公司将股票首次向公众投资者出售。投资银行负责协助组织此次发行，协调信息披露和营销工作，评估投资者的需求，并支持确定发行价格和分配股票。

在IPO中，可能会包含普通股、优先股，或者两者都有。普通股是由公司新发行的股份，为发行人提供新的资金。而优先股则由现有股东出售，从而将所有权转移给购买者，同时不会为公司筹集新的资金。因此，发行股票的构成会影响到公司的融资效果以及销售所得资金的分配情况。

在IPO定价时存在困难，因为该股票的公开交易价格尚未确定。估值过程可能会基于财务预测、同类公司的状况、以往的交易情况以及贴现现金流分析来进行。同时，投资者的反馈也能提供有关需求的额外信息。这些输入数据通常用于确定一个估值范围，而不是一个单一的可观测数值。

在存在这种不确定性的情况下，发行价格需要慎重制定。较高的价格虽然能增加每股的收益，但可能会降低需求；而较低的价格则有利于股份的分销，但如果股份后的价格上升，那么初始投资者可能会获得更大的收益。因此，定价决策必须综合考虑发行公司的资金状况以及成功发行所需的条件。

IPO~抑价通常指从发行价到首日市场价格的正初始回报，常用首日收盘价计量。其幅度因企业、时期和市场而异，可能反映信息摩擦、不确定性、投资者参与激励或发行流程特征。

\begin{tcolorbox}[
  enhanced,
  breakable,
  colback=white,
  colframe=FinanceBlue,
  boxrule=0.7pt,
  sharp corners,
  borderline west={2.2pt}{0pt}{FinanceBlue},
  title=\textbf{财务展望：盘后定价不足是一种结果现象},
  coltitle=white,
  colbacktitle=FinanceBlue,
  fonttitle=\bfseries,
  before skip=1.2em,
  after skip=1.2em
]

在二级市场交易开始之后，才会出现首日上涨的情况。相比之下，发行价格是在对投资者需求以及股票交易价格不确定的情况下预先确定的。

因此，抑价本身并不表明发行人或承销商能够事先知道后来的市场价格。这是一项事后结果，可能同时反映定价决策，以及交易开始时得到解决的不确定性。

\end{tcolorbox}

\subsection{债券发行与定价}

承销工作主要围绕确定债券向投资者发行的条件展开。市场收益率、发行人的信用状况、到期日、清偿顺位、流动性以及合同条款都影响着投资者所需的回报水平。

因此，定价始于对市场的参考点。政府的票息曲线、二级市场的信用利差、近期发行情况以及可比的借款人提供了所需票息的基准。承销商和发行人利用这些基准数据，结合投资者需求的迹象，来确定合适的票息和发行价格。

对于固定利率的债券，票面利率不必等于市场要求的收益率。当票息与投资者要求的收益率不同时，债券可以以高于或低于面值的价格发行。实际上，票息和发行价格通常会一起确定，这样债券就可以与类似的证券进行直接比较，并且可以以合理的条件进行分配。

\begin{example}{从市场基准到最终发行条款}{bond-issue-pricing}

有一家公司计划发行为期五年的债券。该政府债券的收益率为2.8\%，而同类企业债券的利率则比政府债券高出约120个基点。

与投资者的初步讨论表明，对债券的需求在110个基点范围内较为有限，而在125个基点附近则存在较强的需求。考虑到这些迹象，发行人和承销商最终决定将发行利率定为125个基点，这意味着债券的发行收益率大约为

\[
2.8\%+1.25\%=4.05\%.
\]

票息和发行价格可以围绕这个要求的收益率来设定。示例展示了市场基准和投资者反馈如何共同影响新债券发行的最终条款。

\end{example}

投资者的需求可能因期限、信用质量、货币、流动性以及其他特征而有所不同。因此，来自银行、保险公司、养老基金、资产管理公司以及其他投资者的反馈可能会影响债券的最终期限、发行规模、定价以及资金分配。

债务发行的条件可能会继续调整，直到发行方的融资目标和投资者的必要报酬率足够接近，使得该交易能够顺利进行。

\begin{figure}[H]

\centering

\begin{tikzpicture}[
  >=Latex,
  box/.style={
    draw,
    rounded corners,
    align=center,
    minimum width=3.1cm,
    minimum height=0.9cm,
    inner sep=5pt
  },
  flow/.style={-Latex, line width=0.85pt, FinanceBlue}
]

\node[box] (bench) at (-4.6,1.25)
{市场基准数据\\和信贷条件};

\node[box] (demand) at (-4.6,-0.75)
{投资者需求\\和偏好};

\node[box] (yield) at (0,0.25)
{所需收益率};

\node[box] (terms) at (4.0,0.25)
{票息\\发行价格};

% Two inputs converge into a common vertical line
\coordinate (mergeTop) at (-2.0,1.25);
\coordinate (mergeBottom) at (-2.0,-0.75);
\coordinate (mergeCenter) at (-2.0,0.25);

\draw[line width=0.85pt, FinanceBlue] (bench.east) -- (mergeTop);
\draw[line width=0.85pt, FinanceBlue] (demand.east) -- (mergeBottom);

% Shared connection and output
\draw[line width=0.85pt, FinanceBlue]
  (mergeTop) -- (mergeBottom);

\draw[flow] (mergeCenter) -- (yield.west);
\draw[flow] (yield.east) -- (terms.west);
\end{tikzpicture}

\caption{新发行债券的定价}

\label{fig:ch13-bond-issue-pricing}

\end{figure}

\subsection{簿记建档、定价与投资者配售}

簿记建档是承销商在最终定价和配售前收集投资者需求意向的过程。投资者可表明不同价格或收益率下愿意购买的数量，使承销商评估需求的强度及分布。

让 $q_i(p)$ 表示投资者 $i$ 在价格 $p$ 下的指示需求。总体指示需求为……

\begin{equation}
Q(p)=\sum_{i=1}^{n} q_i(p).
\label{eq:ch13-bookbuilding-demand}
\end{equation}

订单簿提供的信息比单纯的总需求更为丰富。承销商可以观察到订单的规模和价格敏感性、参与投资者的构成，以及随着拟议条款的修改而变化的需求情况。这些信息有助于评估发行的价格以及该需求的不确定性程度。

因此，最终定价是依据订单簿来进行的，而不是遵循纯粹的机械规则。对发行人更为有利的定价可能会提高每证券的收益，但在交易完成之前，需求的变化空间就会变小。在确定最终条款时，发行人和承销商还可能考虑执行风险、投资者参与度以及预期的二级市场状况。

当意向需求与发行数量不同时，配售决定证券如何分配。有些发行主要按比例配售，另一些则允许更大裁量。配售可以考虑订单规模、价格敏感性、投资者特征，以及簿记建档中传递的信息。裁量可能改善发行配售，但若受与交易无关的关系或目标影响，也会产生利益冲突。

\begin{example}{债券簿记建档中的价格敏感性}{bond-bookbuilding-sensitivity}

一家公司计划发行约10亿元人民币的五年期债券。在簿记建档过程中，投资者提交了总计约16亿元人民币的订单，预期的收益率为3.0\%。不过，其中约5亿元人民币的订单是以至少3.2\%的收益率为条件的。

如果最终定价定为3.0\%，那么这些对价格敏感的订单可能就不会被执行，这样一来，与最终定价条款相一致的订单就只剩下约11亿元人民币的价值了。不过，随着市场情况的变化，这些订单也可能在分配之前被重新调整。

因此，承销商必须评估提交订单的总量，以及它们的价格敏感性和可靠性。设定较高的收益率可能会吸引更多的需求，但这也会增加发行人的借款成本。因此，簿记建档涉及到平衡投资者的需求与发行人面临的融资成本之间的关系。
\end{example}

\section[做市与二级市场中介]{做市与二级市场中介}

在二级市场上，交易商通过报价、执行客户订单以及当买卖订单不一致时临时持有证券来支持交易活动。这些交易活动利用资产负债表的能力来吸收短期的库存不平衡问题，并促进可交易证券的交易。

\subsection{交易商库存、资产负债表与流动性供给}

交易商可以在找到其他愿意购买该证券的交易商之前，从客户那里购买该证券。这笔交易会进入交易商的库存中，这样客户就可以进行交易，而无需等待相反的订单。

令 $I_t$ 表示周期 $t$ 结束时交易商的库存，$Q_t^{B}$ 表示从客户处采购的商品数量，$Q_t^{S}$ 则表示向客户销售的商品数量。一个简单的库存关系可以表示为：

\begin{equation}
I_t=I_{t-1}+Q_t^{B}-Q_t^{S}.
\label{eq:ch13-dealer-inventory}
\end{equation}

持有库存会使交易商面临价格波动的风险，同时还需要资金支持和资产负债表能力。如果交易商从客户那里购买债券，而在债券出售之前市场收益率上升，那么库存的价值就会下降。因此，资本需求、融资成本、套期保值成本以及内部风险限制等因素都会影响到交易商接受客户订单的意愿。

\begin{example}{客户交易中的库存风险}{dealer-inventory-risk}

该交易商购买了一份市场价值达5000万元人民币的债券头寸，并计划当天晚些时候将其转售。如果在出售该头寸之前，市场价值下降0.4\%，那么其实际损失将约为……

\[
50{,}000{,}000\times 0.004=\text{RMB }200{,}000.
\]

\end{example}

因此，交易商的流动性供应会随着市场状况的变化而不同。当波动性上升时，融资成本会增加，或者风险限制会变得更加严格，交易商可能会持有较少的库存或给出的报价不那么积极。即使证券的合同条款没有发生变化，流动性也可能减弱。

承担库存的作用，使交易商中介不同于简单订单撮合。交易平台可以不持有头寸而连接买卖双方，交易商则通过暂时使用自身资产负债表提供即时成交服务。

\subsection[买卖价差、信息与交易成本]
{买卖价差、信息与交易成本}

交易商通常会报出一个买入价，表示它愿意购买；同时也会报出一个卖出价，表示它愿意出售。所报的价差就是……

\begin{equation}
s=P_{\mathrm{ask}}-P_{\mathrm{bid}}.
\label{eq:ch13-dealer-spread}
\end{equation}

这种价差反映了多种交易成本和风险来源。订单处理和市场准入涉及直接成本，而库存头寸则使交易商面临价格和资金风险。逆向选择增加了信息因素，因为当交易商与那些对证券价值有更好了解的投资人进行交易时，就会引入更多的信息因素。

这些证券的大小因证券和交易的不同而有所差异。对于流动性较高的证券来说，小额交易可能不需要大量的库存投入；而对于流动性较低的证券而言，大额交易则更难以进行对冲处理，从而让交易商面临更大的价格风险。

\begin{tcolorbox}[
  enhanced,
  breakable,
  colback=white,
  colframe=FinanceBlue,
  boxrule=0.7pt,
  sharp corners,
  borderline west={2.2pt}{0pt}{FinanceBlue},
  title=\textbf{财务展望：所报价差并非全部交易成本的体现},
  coltitle=white,
  colbacktitle=FinanceBlue,
  fonttitle=\bfseries,
  before skip=1.2em,
  after skip=1.2em
]

报价买卖价差描述给定数量下可得的价格。如果最优报价提供的数量有限，较大交易可能需要按较不利的价格执行。

因此，交易成本既可能包括报价价差，也可能包括执行所需数量时的额外价格让步。交易规模相对于市场可得流动性较大时，这一区别更加重要。

\end{tcolorbox}

信息不对称还会进一步影响交易商的报价。如果接到的订单中包含更多关于未来价值的信息，那么针对这些订单进行交易可能会给交易商带来损失。更宽的价差、更小的报价数量或者不那么激进的定价可以降低这种风险，不过这也会增加其他投资者即时执行交易的成本。

因此，买卖价差反映了在不确定性情况下提供即时交易服务的成本。这些价差往往取决于证券的流动性、交易的规模、交易商的风险敞口，以及交易发生的信息环境。

\section{顾问服务与战略交易}

投资银行还为企业提供关于那些涉及所有权、控制权或资本结构变化的交易的建议。他们的职责是帮助评估各种选择方案，协商交易条款，在需要时安排融资，并协调那些其价值依赖于企业特定信息和合同设计的交易的执行过程。

\subsection{并购顾问}

并购活动涉及对企业资产和现金流的控制权的转移或重组。投资银行可以通过评估交易方案、估算价值、识别交易对手方，以及支持谈判和执行过程，为收购方、出售方、目标企业或专门委员会提供咨询建议。

交易决策不仅仅涉及估值问题。买家可能会期望通过合并两家公司来获得运营或财务上的好处，但能否实现这些好处仍然存在不确定性。卖方在考虑价格的同时，还需要考虑支付方式、融资条件、监管批准情况以及交易能够成功完成的概率等因素。因此，估值为谈判提供了依据，而并非是决定交易价格的唯一依据。

一个简单的框架将独立价值与从组合中获得的增量价值区分开来。令 $V_{acquirer}$ 和 $V_{target}$ 分别表示收购方和目标公司的独立价值，$S$ 表示将它们组合后预期的总收益，$C$ 则表示交易和整合成本的现值。组合后的价值就是……

\begin{equation}
V_{merger}=V_{acquirer}+V_{target}+S-C.
\label{eq:ch13-combined-transaction-value}
\end{equation}

如果收购方以总对价$P_{target}$收购目标企业，那么在这一简化框架下，收购方现有股东预期的净现值将为……

\begin{equation}
\textit{NPV}_{acquirer}=S-C-(P_{target}-V_{target}).
\label{eq:ch13-acquirer-transaction-npv}
\end{equation}

术语“$P_{target}-V_{target}$”指的是相对于目标资产预估独立价值的收购溢价。此类交易可能会带来正的合并收益，而这些收益中的大部分将通过协商确定的溢价分配给目标股东。

\begin{example}{协同效应与协商溢价}{ma-synergy-premium}

该目标项目的预计独立价值约为10亿元人民币。买家期望获得约2.5亿元人民币的总收益，而整合成本现值约为5000万元人民币。因此，预期净增量价值为2亿元人民币。

如果约定的购买价格为11.5亿元人民币，那么相对于目标资产预估独立价值的溢价约为1.5亿元人民币。基于这些假设，买方预期的交易净现值如下：

\[
200-150=\text{RMB }50\text{ 百万}.
\]

\end{example}

顾问帮助客户评估这些权衡因素，并将估值结果转化为具体的交易条款。咨询服务可以提高决策的质量和执行效率，但它并不能确保该交易最终能够创造价值。

\subsection{重组、资本重整与困境交易}

企业可能会出于战略原因或应对财务压力而调整其资本结构。重组会改变金融请求权的组合、期限或所有权结构，而不一定改变企业的运营资产。当支付义务变得难以维持时，重组可能需要对债务条款、优先级、回收权或控制权进行更实质性的调整。

投资银行可以为公司、债权人或投资者提供关于这些替代方案的评估建议。在重组过程中，分析可能会关注杠杆率、债务结构、融资灵活性以及权益和控制的分配。在财务困境的情况下，更多的注意力会集中在企业价值、预期回收情况、短期融资需求以及各索赔群体之间的价值分配上。可能的交易包括债务或权益的发行、债务期限延长、交换提议、新的融资方式以及资产出售。

一种负债管理交易体现了这一机制。面临近期到期义务的公司可能会提供期限较长的证券作为抵押品来偿还现有债务。延长到期时间可以降低即时的再融资压力，但债权人会对比新提供的证券与原有证券的价值、优先级、抵押品以及预期的回收金额。因此，这些条款必须能够为足够多的持有人提供足够的价值，以吸引他们的参与。

回收行为也可以重新分配价值以及控制权。高级债权人、次级债权人和股东在回收方式、风险以及对公司未来控制权的看法上可能有所不同。顾问可以协助进行估值、交易设计以及谈判，而每个债权人的可执行权利则仍然受合同和适用法律的约束。

\subsection{估值、谈判与交易执行}

估值提供了一个比较交易方案的工具，但它通常依赖于假设，而不是得出一个唯一确定的价值。贴现现金流分析依赖于预测和贴现率，而同类公司分析和先例交易分析方法则取决于所选基准指标的相关性。当预期的协同效应、融资能力或控制优势不同时，战略价值也可能因购买方的不同而有所差异。因此，投资银行通常会使用估值范围和敏感性分析来展示估计价值如何随着关键假设的变化而变化。

谈判过程会将这些评估结果转化为具体的交易条款。价格只是其中一个因素而已。支付方式可以是现金、股票、债务，或者多种方式的组合；而绩效补偿条款、条件性支付、退出费、融资条件以及结算条款则会影响风险和价值的分配。因此，虽然这些交易的价格看起来相似，但实际上它们给参与各方带来的经济负担却可能大相径庭。

执行环节决定了是否能在约定的条款下完成所协商的交易。顾问们负责协调融资、文件编制、时间表安排以及信息流通，这些工作涉及法律、会计、监管等多个领域。融资条件的变化、监管审查、股东批准、尽职调查的结果或市场状况都可能影响交易的完成概率和条款。

\begin{tcolorbox}[
  enhanced,
  breakable,
  colback=white,
  colframe=FinanceBlue,
  boxrule=0.7pt,
  sharp corners,
  borderline west={2.2pt}{0pt}{FinanceBlue},
  title=\textbf{财务展望：价值和价格并非等同的概念。},
  coltitle=white,
  colbacktitle=FinanceBlue,
  fonttitle=\bfseries,
  before skip=1.2em,
  after skip=1.2em
]

估值是指在既定假设条件下，对某项交易的经济价值进行估算。协商达成的价格不仅反映了这些估值结果，还考虑了谈判条件、竞争报价、融资条款、合同保护以及执行风险等因素。

因此，定价范围可以为谈判提供指导，但并不能直接确定最终的交易价格。

\end{tcolorbox}

%========================================================
% Summary
%========================================================

\TailModule{summary}{本章小结}

投资银行通过信息生产、认证、证券设计、分销与执行，支持资本市场交易。其作用不同于传统资产负债表中介，因为投资者通常仍持有融资实体直接发行的请求权。

承销有助于将证券从发行人转移给投资者。公司的承诺使承销商面临配售风险，同时承销还涉及到股份分配的能力和风险，而簿记建档则提供了定价和分配的信息。

在二级市场中，交易商利用库存和资产负债表能力来提供即时性。买卖价差反映了交易成本、库存风险、信息风险以及当前的市场状况。

投资银行还为并购、资本重组和重组等业务提供咨询。在这些活动中，它们的经济作用主要体现在信息收集、声誉建设、估值分析、交易专长以及合理运用资产负债表能力等方面。

%========================================================
% Key Terms
%========================================================

\TailModule{keyterms}{关键术语}

投资银行（Investment bank）；资本市场中介（capital-market intermediation）；认证（certification）；尽职调查（due diligence）；证券设计（security design）；承销（underwriting）；代销（best efforts）；承销价差（underwriting spread）；承销团（underwriting syndicate）；首次公开发行（initial public offering）；公司新发行股份（primary shares）；现有股东出售股份（secondary shares）；债券承销（debt underwriting）；簿记建档（bookbuilding）；认购意向（indication of interest）；投资者配售（investor allocation）；做市（market making）；交易商库存（dealer inventory）；买卖价差（bid--ask spread）；战略顾问（strategic advisory）；合并（merger）；收购（acquisition）；收购溢价（acquisition premium）；资本重整（recapitalization）；重组（restructuring）；负债管理（liability management）；交易执行（transaction execution）。

%========================================================
% Financial AI Lab
%========================================================

\TailModule{aiprompt}{金融~AI~实验室}

\textbf{实验室：对人工智能融资建议的审核}

请选择一个最近已筹集到权益或债务资本的上市公司。请一个AI系统担任投资银行顾问，并推荐一种发行结构。该建议应明确融资目标、证券类型、发行规模、可能的投资者群体、相关的估值或收益率基准，以及主要的执行风险。

将建议分为四个阶段进行审核。

\begin{enumerate}

\item \textbf{确认融资需求以及现有的资本结构。}

利用财务报表和正式披露信息来确认公司的债务、权益、流动性状况，以及资金的使用方式。将经过验证的事实与人工智能所提出的假设区分开来。

\item \textbf{评估证券设计。}

评估所提出的证券条款是否符合发行人的融资需求以及潜在的投资者群体。对于债务融资，需要考虑到期日、还款结构、清偿顺位以及相关的抵押品。对于权益融资，则需要考虑发行规模、所有权效应以及发行形式。确定一种替代的融资结构，并解释其中的主要权衡。

\item \textbf{审计定价证据。}

对于权益发行，将拟定的估值与可观察到的市场价格或合适的可比公司数据进行比较。对于债券发行，将拟定的收益率或息差与适当的基准曲线以及可比发行人进行比较。确保日期、货币、到期日以及单位都是一致的。

\item \textbf{与承销承诺分开考虑的建议。}

确定推荐意见中的哪些部分涉及信息、估值和交易建议，而哪些部分可能会让投资银行面临承销风险。解释在 firm commitment 模式下与 placement 或 best-efforts 安排相比，分配风险是如何变化的。

\end{enumerate}

%========================================================
% Discussion Questions
%========================================================

\TailModule{questions}{讨论题}
\begin{enumerate}

\item 为什么在证券发行中使用投资银行并不能直接将融资转化为间接融资呢？

\item 投资银行声誉如何能够减少信息摩擦？为什么认证并不能消除投资者判断的重要性呢？

\item 一家承诺型承销商承担什么样的分销风险？而一般的代理公司通常并不承担这种风险。

\item 为什么在BookBuild交易中，最终确定的报价可能会与某些投资者报价的最高价格有所不同呢？

\item 交易商如何支持流动性供应，以及交易商承担哪些资产负债表风险？

\item 在收购过程中，为什么目标公司的股东能够获得巨大的收益，而收购方现有股东的资产价值却几乎为零呢？

\end{enumerate}

%========================================================
% Applied Problems
%========================================================

\TailModule{appprobs}{应用题}

\begin{enumerate}

\item \textbf{承销经济学。}

一家公司在公开募股中发行了3000万股股票，公开价格约为每股18元人民币。承销团以每股17.55元人民币的价格从发行人手中购得了这些股票。

\begin{enumerate}
\item 计算每股的承销总差额。
\item 计算总毛差幅。
\item 计算发行人在支付其他费用之前所收到的收益。
\item 解释为什么总利差与承销团队的利润并不相同。
\end{enumerate}

\item \textbf{主要和次要份额。}

此次公开募股中，公司发行了8000万股优先股，同时现有股东还出售了2000万股普通股。本次发行的价格约为每股12元人民币。

\begin{enumerate}
\item 计算与优先股相关的总收入。
\item 计算与次级股份相关的总收入。
\item 在扣除费用和开支之后，公司还能获得多少资金呢？
\item 解释这两个组成部分如何不同地影响所有权和融资问题。
\end{enumerate}

\item \textbf{簿记建档与需求。}

发行方计划出售5000万股股票。预计的需求量为：以人民币20元价格交易的股票为9000万股，以人民币21元价格交易的股票为6500万股，以人民币22元价格交易的股票为4200万股。

\begin{enumerate}
\item 在哪些定价条件下，需求会超过供应的数量呢？
\item 为什么仅这些信息还不足以确定一个合适的报价呢？
\item 确定订单簿中承销商除了总体数量之外还可以审查的两种定性特征。
\item 解释更高的发行价格如何同时影响发行公司的收益以及执行风险。
\end{enumerate}

\item \textbf{交易商库存。}

该交易商在一天的开始时持有市值约为1000万人民币的债券。它从客户那里购买了市值约为3500万人民币的债券，并在当天出售了价值2800万人民币的债券。

\begin{enumerate}

\item 计算最终的库存数量，以所述交易价值为准。

\item 如果最终库存在出售前价值下降了0.6\%，那么请计算出相应的市价损失。

\item 请列出另外两项可能影响提供流动性的经济因素。

\item 解释更严格的内部风险限制如何即使债券的合同现金流保持不变，也能影响客户的执行。

\end{enumerate}

\item \textbf{收购价格与谈判事宜。}

该目标项目的预估独立价值为20亿元人民币左右。投标方估计的 총协同效应为4.2亿元人民币，而整合和交易成本现值总计为1.2亿元人民币。

\begin{enumerate}
\item 计算预计的净增量价值。
\item 如果投标者提供的报价为22.4亿元人民币，那么相对于单独项目的价值而言，这一报价所包含的溢价是多少呢？
\item 使用公式~(\ref{eq:ch13-acquirer-transaction-npv})，在给定假设条件下计算投标方预期的交易净现值。
\item 如果整合成本上升到2.2亿元人民币，请重新计算净现值。
\item 解释为什么协商出的价格可能与基于这些假设的估值结果有所不同。
\end{enumerate}

\item \textbf{债券发行与市场反馈。}

有一家公司计划发行为期五年的债券。该政府债券的收益率为3.0\%。相比之下，同类企业债券的收益率比政府债券高出约110个基点。初步的投资者反馈显示，在100个基点的价差下，市场需求较弱；而在120个基点的价差下，市场需求则较为旺盛。

\begin{enumerate}

\item 计算在110个基点利差下所预期的基准收益率。

\item 计算在120个基点利差下所暗示的发行收益率。

\item 请解释为什么发行方即使接受更高的利差也是可以接受的，尽管这会增加借款成本。

\item 在定价之前，除了基准收益率之外，还有两个因素可能导致最终还款条款发生变化。

\end{enumerate}

\end{enumerate}

%========================================================
% Case Analysis
%========================================================

\TailModule{caseanalysis}{案例分析}

\textbf{公开发行}

阿里巴巴集团控股有限公司于~2014~年~9~月在美国首次公开发行，定价为每份美国存托股（ADS）68~美元。基础发行包括~320,106,100~份~ADS：其中阿里巴巴发行~123,076,931~份新增~ADS，现有股东出售~197,029,169~份~ADS。按发行价格计算，基础发行总所得约为~217.7~亿美元。

六家投资银行担任联合簿记管理人。招股说明书披露，每份~ADS~的承销折扣和佣金为~0.816~美元。阿里巴巴及部分出售股份的股东还授予承销商选择权，最多可额外购买~48,015,900~份~ADS。该选择权随后被全部行使，使总发行所得增至约~250.3~亿美元。

运用该交易考察投资银行的多项功能。首先，区分阿里巴巴出售的新增~ADS~与现有股东出售的存量~ADS。哪部分为公司筹集新资本？哪部分将所有权从原股东转移给新投资者？

其次，计算承销折扣占~68~美元发行价的百分比。这项报酬反映哪些服务和风险？为什么不应将其解释为纯承销利润？

第三，考察联合簿记管理人的作用。承销团如何扩大投资者覆盖、在簿记建档中汇集信息，并分担分销风险暴露？发行人希望提高发行价格，与承销商关注执行风险和投资者参与之间，可能产生什么张力？

最后，解释额外购买~ADS~的选择权，如何在基础发行规模之外提供灵活性并支持分销过程。分析应聚焦承销与分销，而非此后股价表现。

\emph{主要来源：阿里巴巴集团控股有限公司，于2014年9月向美国证券交易委员会提交的最终招股说明书；阿里巴巴集团关于IPO定价以及承销商选择权的公告，2014年9月。}

%========================================================
% Further Reading
%========================================================

\TailModule{furtherread}{延伸阅读}

\begin{itemize}

\item Beatty, Randolph P., and Jay R. Ritter (1986),
``Investment Banking, Reputation, and the Underpricing of Initial Public Offerings（投资银行、声誉与首次公开发行抑价）,''
\emph{Journal of Financial Economics}, 15(1--2), 213--232.

\item Benveniste, Lawrence M., and Paul A. Spindt (1989),
``How Investment Bankers Determine the Offer Price and Allocation of New Issues（投资银行家如何决定新发行证券的发行价格与配售）,''
\emph{Journal of Financial Economics}, 24(2), 343--361.

\item Chemmanur, Thomas J., and Paolo Fulghieri (1994),
``Investment Bank Reputation, Information Production, and Financial Intermediation（投资银行声誉、信息生产与金融中介）,''
\emph{The Journal of Finance}, 49(1), 57--79.

\item Ellis, Katrina, Roni Michaely, and Maureen O'Hara (2000),
``When the Underwriter Is the Market Maker: An Examination of Trading in the IPO Aftermarket（当承销商兼任做市商：对IPO发行后市场交易的考察）,''
\emph{The Journal of Finance}, 55(3), 1039--1074.

\item Servaes, Henri, and Marc Zenner (1996),
``The Role of Investment Banks in Acquisitions（投资银行在收购中的作用）,''
\emph{The Review of Financial Studies}, 9(3), 787--815.

\end{itemize}

\ChapterTail{summary,keyterms,aiprompt,questions,appprobs,caseanalysis,furtherread}
